\mode*
\section{Semantic markup}
\label{sec:markup}

The third educational objective is markup as \emph{meaning}: sections,
emphasis, quotations, cross-references and citations declare what things
\emph{are}; how they \emph{look} is decided elsewhere.

\subsection{Documented difficulties}

The seed round found no studies of learners' understanding of semantic
versus presentational markup.
% TODO(#1): The systematic phase should look at markup/HTML education studies
% and end-user document structuring; course data (datintro LaTeX lab
% grading, which specifically checks semantic-markup use) supplies the
% local observations.

\subsection{Candidate critical aspects}

% XXX Preliminary.
\begin{itemize}
  \item \emph{Say what it is, not how it looks}: \verb|\section| versus
    large bold text; \verb|\emph| versus italics.
  \item \emph{One declaration, many appearances}: the same source
    renders differently under a different class or style.
  \item \emph{References are relations}: \verb|\ref|/\verb|\cite| link to
    things, and the compiler maintains the links.
\end{itemize}

\subsection{Patterns of variation}

\begin{itemize}
  \item Contrast (markup varies, appearance invariant): a hand-formatted
    \enquote{heading} (bold, large) next to \verb|\section|---identical
    look, then observe the table of contents and numbering diverge.
  \item Contrast (style varies, markup invariant): compile the same
    semantically marked-up source under two document classes.
  \item Fusion: renumber by inserting a section early in the
    document---automatic numbering, references and table of contents all
    update.
\end{itemize}
